\section{Uniqueness after canonical normalization}\label{sec:uniqueness}

We now remove the offset parameter.  The preceding section gives a
continuous family $G_s$ of canonically normalized hierarchies for
$s\in I=(-R,R)$, together with its positive second-moment length
$\ell(s)$ before normalization.  Two descriptions of the same microscopic
model relate different members of this family.  An additional ordinary
step changes the physical mesh by $L^{-1}$; an inclined first step changes
it by $5^{-1}$ and rotates it.  The two reflected inclinations have the
same offset shift.  These facts supply both rotational invariance and two
incommensurable periods in $s$.
Figure~\ref{fig:two-scales} records the two descriptions and their
physical length factors.

For a rotation $U$ of the Euclidean plane, write $U_*G$ for the hierarchy obtained by
pushing each insertion coordinate forward by $U$.  At the field level
this means $(U_*\phi)(f)=\phi(f\circ U)$.  Component indices are unchanged:
these are rotations of Euclidean coordinates, distinct from the internal $O(3)$ symmetry.

\subsection{Two descriptions of one cutoff}

\begin{proposition}[Depth and inclination relations]\label{prop:two-scales}
For $s,s-\gamma\in I$,
\begin{equation}\label{eq:uni-depth}
 G_{s-\gamma}=G_s,
 \qquad \ell(s-\gamma)=L^{-1}\ell(s).
\end{equation}
Let $d=A_+-A_{\mathrm r}=A_--A_{\mathrm r}$ be the offset shift
defined in Section~\ref{sec:trajectories}.  For $s,s+d\in I$,
\begin{equation}\label{eq:uni-incline}
 \ell(s+d)=5^{-1}\ell(s),\qquad
 G_{s+d}=(O_+^{-1})_*G_s=(O_-^{-1})_*G_s.
\end{equation}
\end{proposition}

\begin{proof}
Fix $s,s-\gamma\in I$ and take bare couplings
$\beta_N=\vartheta_N+s$.  Regard the same microscopic law first as an
ordinary trajectory of depth $N$, and then as one of depth $N+1$.
The latter offset is
\[
 \beta_N-\vartheta_{N+1}
 =s-(\vartheta_{N+1}-\vartheta_N)\longrightarrow s-\gamma.
\]
Both descriptions are admitted for all sufficiently large $N$ by
Proposition~\ref{prop:admission}.  At every cutoff, canonical
normalization in Proposition~\ref{prop:normalization} yields the same
field $\Psi_{\beta_N}$.  Their limiting hierarchies are therefore equal.
Their cutoff second-moment lengths before normalization are respectively
$L^{-N}\xi_{\beta_N}$ and $L^{-(N+1)}\xi_{\beta_N}$.  Taking their
positive limits proves \eqref{eq:uni-depth}.

Next use an ordinary and an inclined trajectory of the same depth $N$
and the same bare coupling $\vartheta_N+s$.  By
Proposition~\ref{prop:offset-limits}, the inclined limiting hierarchy
is the ordinary one at offset $s+d$.  Its microscopic spacing is
$L^{-N}/5$ and its spatial embedding is $O_\pm^{-1}$.
Thus its second-moment length is exactly one fifth of the ordinary
one, while its canonically normalized field is the pushforward of
$\Psi_{\beta_N}$ by $O_\pm^{-1}$.  Passing to the limit gives
\eqref{eq:uni-incline}.  Reflection of the microscopic construction
identifies the scalar canonical increments of the two inclinations,
which is why the same $d$ occurs in both equalities.
\end{proof}

\begin{figure}[tb]
\centering
\begin{tikzpicture}[>=Latex,node distance=14mm and 24mm,
 every node/.style={font=\small,align=center},
 box/.style={draw=MidnightBlue!65,rounded corners,fill=MidnightBlue!3,
             inner sep=6pt,text width=38mm}]
\node[box] (a) {Ordinary, depth $N$\\
 $s,\quad a=L^{-N}$\\$G_s,\quad\ell(s)$};
\node[box,right=of a] (b) {Ordinary, depth $N+1$\\
 $s-\gamma+o(1)$\\$G_{s-\gamma},\quad\ell(s)/L$};
\node[box,below=of a] (c) {Inclined, depth $N$\\
 $s,\quad a=L^{-N}/5$\\$G_{s+d},\quad\ell(s)/5$};
\node[right=of c,text width=38mm] (d) {All three use the same\\bare coupling $\beta_N$.};
\draw[->,thick] (a)--node[above] {one extra step} (b);
\draw[->,thick] (a)--node[left,text width=22mm] {inclined\\first step} (c);
\draw[->,MidnightBlue!70] (d.north)--(b.south);
\draw[->,MidnightBlue!70] (d.west)--(c.east);
\end{tikzpicture}
\caption{The two scale comparisons.  The displayed offsets in the boxes
are the bare offsets of the trajectories; the inclined hierarchy matches
the ordinary hierarchy at $s+d$.  Canonical normalization cancels their
different field multipliers.  The inclined comparison retains only the
Euclidean rotation, as in \eqref{eq:uni-incline}.}
\label{fig:two-scales}
\end{figure}

\subsection{An elementary uniqueness mechanism}

We isolate the topological part of the argument.  Its input consists only
of continuous distributions, positive lengths, and the relations just
proved.  In particular, it requires no monotonicity of the shooting map.

\begin{lemma}[Two scale relations force constancy]\label{lem:dense-periods}
Let $L>1$ be a power of two, let $\gamma>0$, and let
$I=(-R,R)$.  Suppose $\ell:I\to(0,\infty)$ is continuous and
$s\mapsto G_s$ is a family of distribution hierarchies whose evaluation
on every Schwartz test is continuous.  Let $O_\pm$ be rotations by
$\pm\vartheta$, where $\cos\vartheta=3/5$ and
$\sin\vartheta=4/5$.  Suppose \eqref{eq:uni-depth} and
\eqref{eq:uni-incline} hold wherever both offsets belong to $I$, and
suppose $2R-|d|>\gamma$.
Then every $G_s$ is Euclidean rotation invariant, and $G_s$ is independent
of $s$.
\end{lemma}

\begin{proof}
First take $s$ in the interval $J=I\cap(I-d)$.  Equality of the two
pushforwards in \eqref{eq:uni-incline} implies invariance of $G_s$
under rotation by $2\vartheta$.  The number $\vartheta/\pi$ is
irrational.  Indeed, if it were rational, $e^{i\vartheta}$ would be a
root of unity, making
$e^{i\vartheta}+e^{-i\vartheta}=6/5$ an algebraic integer.  A rational
algebraic integer is an integer, which $6/5$ is not.  Powers of this
rotation are consequently dense in $SO(2)$.  The action of rotations
on Schwartz tests is continuous, so every distribution in $G_s$ is
invariant under all Euclidean rotations.  Equation~\eqref{eq:uni-incline}
therefore gives $G_{s+d}=G_s$ on $J$.

To use these local relations globally, extend the two families to
$\R$.  For $x\in\R$, choose an integer $n$ with $x+n\gamma\in I$ and
set
\begin{equation}\label{eq:uni-extension}
 \widetilde G_x=G_{x+n\gamma},\qquad
 \widetilde\ell(x)=L^{-n}\ell(x+n\gamma).
\end{equation}
Such an $n$ exists because $|I|>\gamma$.  These definitions do not
depend on the choice: two representatives can be joined by consecutive
$\gamma$-steps staying in the interval $I$, and
\eqref{eq:uni-depth} applies at every step.  On overlapping translates
of $I$, the definitions agree; hence the extended families are continuous.
They satisfy
\[
 \widetilde G_{x+\gamma}=\widetilde G_x,
 \qquad \widetilde\ell(x+\gamma)=L\widetilde\ell(x).
\]
The interval $J$ has length $2R-|d|>\gamma$, so every class modulo
$\gamma$ has a representative in $J$.  Choosing that representative in
\eqref{eq:uni-extension} transports both the rotation invariance and
the $d$-relations to every $x\in\R$:
\begin{equation}\label{eq:uni-global-periods}
 \widetilde G_{x+d}=\widetilde G_x,
 \qquad \widetilde\ell(x+d)=5^{-1}\widetilde\ell(x).
\end{equation}

If $d/\gamma=p/q$ with integers $p$ and $q>0$, then iterating the
length relations gives
\[
 5^{-q}\widetilde\ell(x)
 =\widetilde\ell(x+qd)
 =\widetilde\ell(x+p\gamma)
 =L^p\widetilde\ell(x).
\]
Since $\widetilde\ell(x)>0$, this would imply $5^{-q}=L^p$,
contradicting unique prime factorization.  Thus $d/\gamma$ is irrational.
The subgroup $\Z\gamma+\Z d$ is dense in $\R$.  Each test evaluation
of $\widetilde G$ is continuous and has this subgroup as periods, and
is therefore constant.  Equality on every test proves equality of the
hierarchies.  Restriction to $I$ proves the claim.
\end{proof}

\subsection{Every bare coupling and the physical conclusions}

\begin{proof}[Proof of Theorem~\ref{thm:canonical}]
Finiteness and strict positivity of $\chi_\beta$ and $\xi_\beta$ for
sufficiently large $\beta$ follow from Lemma~\ref{lem:finite-beta}.
Proposition~\ref{prop:two-scales} supplies the hypotheses of
Lemma~\ref{lem:dense-periods}; our choice of $R$ in
Section~\ref{sec:trajectories} gives the required overlap.  Write $G$
for the common normalized hierarchy.

The reference values $\vartheta_N$ tend to infinity, and
$\vartheta_{N+1}-\vartheta_N\to\gamma>0$.  For every sufficiently large
$\beta$, choose $N=N(\beta)$ so that
\[
 \vartheta_N\le\beta<\vartheta_{N+1}.
\]
Then $N(\beta)\to\infty$, and the offsets
$s(\beta)=\beta-\vartheta_{N(\beta)}$ belong to the fixed compact interval
$[0,2\gamma]\subset I$.  Any sequence $\beta_k\to\infty$ has a
subsequence on which these offsets converge to some $s\in I$.
Propositions~\ref{prop:offset-limits} and~\ref{prop:normalization}
give, along that subsequence, convergence of the canonically normalized
fields to $G$, jointly in law and in every joint moment, and strong
convergence of each component Schwinger distribution in $\mathcal S'$.
Moreover,
\[
 L^{-N(\beta_k)}\xi_{\beta_k}\longrightarrow\ell(s)>0,
\]
so $\xi_{\beta_k}\to\infty$ along the same subsequence.

These conclusions hold for a subsequence of every sequence tending to
infinity, with the same limiting hierarchy.  If any asserted convergence
failed, a sequence staying outside a fixed neighborhood of the claimed
limit would have such a subsequence, a contradiction.  This proves the
full $\beta\to\infty$ assertions.  The same argument applies to a
neighborhood in the strong topology of $\mathcal S'$; it does not require
that topology to be metrizable.  The finite-list laws are consistent, and
their local exponential moments determine them from $G$.

Finally let $\phi^{\mathrm{sh}}$ be the hierarchy in
Theorem~\ref{thm:main}, obtained by fixing a terminal kinetic coupling.
Its auxiliary parameters need not equal those used for the offset
argument.  Lemma~\ref{lem:shooting-normalization} gives finite positive
susceptibility $m_{\mathrm{sh}}$ and second-moment length
$\ell_{\mathrm{sh}}$, and convergence of the corresponding cutoff
quantities.  Thus its normalization
\begin{equation}\label{eq:uni-shooting-identification}
 \phi^{\mathrm{can}}(f)
 =\frac{1}{\ell_{\mathrm{sh}}\sqrt{m_{\mathrm{sh}}}}
    \phi^{\mathrm{sh}}\bigl(f(\,\cdot\,/\ell_{\mathrm{sh}})\bigr)
\end{equation}
is the limit of the same microscopic fields $\Psi_{\beta_N}$ along its
shooting sequence.  Since $\beta_N\to\infty$, this hierarchy is $G$.
Proposition~\ref{prop:normalization} also proves uniqueness of the
positive length and field constants fixing susceptibility and
second-moment length to one.
\end{proof}

\begin{proof}[Proof of Corollary~\ref{cor:canonical-physics}]
Equation~\eqref{eq:uni-shooting-identification} is a positive dilation
of space-time and multiplication of the field by a nonzero constant.
It preserves the Euclidean and internal symmetries, reflection
positivity, symmetry of the Schwinger functions, clustering, and their
factorial distributional bounds.  The reconstruction argument of
Section~\ref{sec:os} therefore applies.

More explicitly, substitution of the rescaled tests and field factors
defines an invertible map of the positive-time polynomial algebras and
preserves their Osterwalder--Schrader inner products.  It induces a
unitary map between the completed reconstructed spaces, taking vacuum
to vacuum.  A time translation by $t$ in canonical coordinates
corresponds to one by $\ell_{\mathrm{sh}}t$ in shooting coordinates.
Consequently the two Hamiltonians satisfy, under this unitary map,
\[
 H_{\mathrm{can}}=\ell_{\mathrm{sh}}H_{\mathrm{phys}}.
\]
The unique vacuum, the nonzero vacuum complement, and the positive
lower bound on the full vacuum-orthogonal Hamiltonian therefore persist.
A connected fourth distribution transforms by the nonzero fourth power
of the field factor and the invertible dilation of its tests.  Strict
separation of time supports is preserved by a positive dilation.
The nonzero separated connected fourth correlation of
Theorem~\ref{thm:main} thus remains nonzero in canonical units.
\end{proof}
